\subsubsection{Fourier 级数}\label{subsec:fourier-series-hilbert-bases}

关于标准正交基的展开最重要的例子是 Fourier 展开.

\begin{Theorem}
\label{thm:fourier-cosine-hilbert-basis}
设 $\phi_0(x)=1$, $\phi_n(x)=\sqrt2\cos(n\pi x)$, $n\in\mathbb{N}$. 那么函数 $\phi_n$, $n=0,1,2,\ldots$, 构成 $L^2(0,1)$ 的一组基.
\end{Theorem}

\begin{Proof}
简单计算表明 $\phi_n$ 为标准正交系. 因而由 Theorem~\ref{thm:hilbert-basis-equivalent-characterizations}, 只需证明 $\phi_n$ 的线性张成空间在 $L^2(0,1)$ 中稠密. 由于 $C([0,1])$ 在 $L^2(0,1)$ 中稠密, 只需证明每个连续函数均可由 $\phi_n$ 的线性组合逼近. 作代换 $\cos\pi x=u$, 它将 $[0,1]$ 双射到 $[-1,1]$. 由 Weierstrass 逼近定理, $[-1,1]$ 上的每个连续函数都可由多项式一致逼近; 因而 $[0,1]$ 上的每个连续函数都可由 $\cos\pi x$ 的多项式一致逼近. 初等三角恒等式表明, 任意表达式 $\sum_{n=0}^N a_n(\cos\pi x)^n$ 都可改写成 $\sum_{n=0}^N b_n\cos(n\pi x)$ 的形式.
\end{Proof}

$L^2(0,1)$ 中的函数也可以展开成正弦级数而不是余弦级数.

\begin{Theorem}
\label{thm:fourier-sine-hilbert-basis}
设 $\psi_n(x)=\sqrt2\sin(n\pi x)$, $n\in\mathbb{N}$. 那么函数 $\psi_n$, $n\in\mathbb{N}$, 构成 $L^2(0,1)$ 的一组基.
\end{Theorem}

\begin{Proof}
利用 $\mathcal{D}(0,1)$ 在 $L^2(0,1)$ 中稠密这一事实. 若 $f\in\mathcal{D}(0,1)$, 则 $f(x)/\sin(\pi x)$ 在 $[0,1]$ 上连续, 由前一定理的证明可知, 它可由形如 $\sum_{n=0}^N a_n\cos(n\pi x)$ 的表达式一致逼近. 因而 $f(x)$ 可由下列形式的表达式一致逼近
\begin{equation}\label{eq:fourier-sine-cosine-product-identity}
\begin{aligned}
\sum_{n=0}^N a_n\cos(n\pi x)\sin(\pi x)
&=\dfrac12\sum_{n=0}^N a_n\bigl(\sin((n+1)\pi x)\\
&\hspace{8em}-\sin((n-1)\pi x)\bigr).
\end{aligned}
\end{equation}
证明完成.
\end{Proof}

Theorems~\ref{thm:fourier-cosine-hilbert-basis} 和~\ref{thm:fourier-sine-hilbert-basis} 给出以下简单推论.

\begin{Corollary}
\label{cor:fourier-complex-hilbert-basis}
函数 $(1/\sqrt2)e^{in\pi x}$, $n\in\mathbb{Z}$, 构成 $L^2(-1,1)$ 的一组基.
\end{Corollary}

\begin{Proof}
$L^2(-1,1)$ 中的任意函数都可以分解为偶部和奇部. 利用前两个定理, 可以将偶部展开成余弦级数, 将奇部展开成正弦级数.
\end{Proof}
